% Generated draft: requires mathematical and source review.
\documentclass[11pt]{article}
\usepackage[margin=25mm]{geometry}
\usepackage{fontspec}
\setmainfont{Latin Modern Roman}
\usepackage{amsmath,amssymb,mathtools,unicode-math}
\setmathfont{Latin Modern Math}
\usepackage{enumitem,needspace,xurl,hyperref}
\setlength{\parindent}{0pt}
\setlength{\parskip}{5pt}
\setlength{\emergencystretch}{3em}
\newcommand{\N}{\mathbb{N}}
\newcommand{\Z}{\mathbb{Z}}
\newcommand{\Q}{\mathbb{Q}}
\newcommand{\R}{\mathbb{R}}
\newcommand{\C}{\mathbb{C}}
\newcommand{\F}{\mathbb{F}}
\newcommand{\A}{\mathbb{A}}
\newcommand{\PP}{\mathbb{P}}
\newcommand{\E}{\mathbb{E}}
\newcommand{\eps}{\varepsilon}
\newcommand{\dd}{\,\mathrm{d}}
\begin{document}
\section*{UnsolvedMath ID: 8800036}
\textbf{Source problem code:} catalogue entry
\subsection{Definitions and mathematical statement}
Let $q$ be a prime power and let $E/\mathbb F_q$ be an elliptic curve. Its rational points form a finite abelian group, which is cyclic if some point generates all of $E(\mathbb F_q)$. Determine how frequently this occurs for a random elliptic curve over $\mathbb F_q$.

The sampling law must be fixed. If $\mathcal E_q$ is the finite set of $\mathbb F_q$-isomorphism classes, one may use the uniform class distribution or the common automorphism-weighted distribution. The latter probability is
\[
P_q^{\mathrm{wt}}=
\frac{\sum_{[E]\in\mathcal E_q,\ E(\mathbb F_q)\ \mathrm{cyclic}}|\operatorname{Aut}_{\mathbb F_q}(E)|^{-1}}
{\sum_{[E]\in\mathcal E_q}|\operatorname{Aut}_{\mathbb F_q}(E)|^{-1}}.
\]
Uniform sampling of nonsingular coefficients in a chosen Weierstrass model is another input convention requiring characteristic-dependent specifications. The source does not select a measure or specify whether an exact formula, an asymptotic, or uniform bounds are requested.

The summary attributes an at-least-$75\%$ lower bound and a strictly-less-than-$100\%$ conclusion to Vladut (1999), but supplies neither the ensemble nor the range/limiting quantifier, so these are not asserted for every field. Its fixed-curve-over-$\mathbb Q$, varying-prime density discussion is a distinct ensemble and includes unstated nonobstruction hypotheses for positive-density claims. It is not substituted for the random-curve question at fixed $q$.
\subsection{Short English statement}
How large is the proportion of elliptic curves over $\mathbb F_q$ with cyclic rational-point group? Fix the random-curve measure and the exact or asymptotic counting regime; fixed-curve prime-density results concern a different question.
\subsection{Sources}
\begin{itemize}
\item Igor Shparlinski, IPAM \emph{Workshop on Number Theory and Cryptography --- Open Problems} (2006), attachment 6224, PDF pages 3 and 18, schedule source \url{https://www.ipam.ucla.edu/programs/workshops/workshop-i-number-theory-and-cryptography-open-problems/?tab=schedule}; accessed 2026-07-29. The supplied provenance records a PDF/legacy-PPT visual audit.
\item Related attributions supplied in the summary: Vladut (1999); Lang--Trotter connections; Cojocaru, Murty, and Duke (2001--2006), with ERH conditions; Bianchi--Tsimerman-type work; Freiberg--Kurlberg--Sober\'on--Vega; and A. Lucchini Arteche. The record supplies no full publication identifiers or precise scope for these references.
\item Local registry \texttt{data/unsolvedmath/problems.json}, permanent UnsolvedMath ID $8800036$, source code catalogue entry, \emph{Frequency of cyclic elliptic-curve groups}. Dataset attribution: Ulam AI and the UnsolvedMath Contributors.
\end{itemize}
\end{document}
